\uuid{H5av}
\titre{Loi uniforme}
\auteur{Maxime Nguyen}
\organisation{AMSCC}
\datecreate{2026-11-24}
\niveau{L2}
\module{Probabilité}
\chapitre{Probabilité continue}
\sousChapitre{Lois usuelles et approximations}
\difficulte{1}
\theme{loi uniforme, densité, espérance, variance}

\contenu{
\begin{question}
Soit $X$ une variable de loi uniforme sur l'intervalle $[2;6]$, notée $X \sim \mathcal{U}([2;6])$. On note $\mathbb{V}(X)$ sa variance et $\mathbf{1}_{[2;6]}$ l'indicatrice de $[2;6]$ ($1$ sur $[2;6]$, $0$ ailleurs). Quelles affirmations sont vraies ?
\begin{answers}
\good{Une densité de $X$ est $f(x) = \frac14\,\mathbf{1}_{[2;6]}(x)$.}
\good{$\mathbb{E}(X) = 4$.}
\bad{$\mathbb{V}(X) = \frac13$.}
\bad{$\mathbb{P}(X = 4) = \frac14$.}
\good{$\mathbb{P}(X \leq 3) = \frac14$.}
\end{answers}
\begin{explanations}
C : pour la loi uniforme sur $[a;b]$, $\mathbb{V}(X) = \frac{(b-a)^2}{12}$, ici $\frac{16}{12} = \frac43$ (le carré a été oublié).

D : $\mathbb{P}(X = 4) = 0$ ; $\frac14$ est la valeur de la densité.

E : $\frac{3 - 2}{6 - 2} = \frac14$.

Contre-sens visé : lire la hauteur de la densité comme la probabilité d'un point.
\end{explanations}
\end{question}
}
